%
% Copyright � 2026 Peeter Joot.  All Rights Reserved.
% Licenced as described in the file LICENSE under the root directory of this GIT repository.
%
%{
\input{../latex/blogpost.tex}
\renewcommand{\basename}{newtons_reciprocal}
%\renewcommand{\dirname}{notes/phy1520/}
\renewcommand{\dirname}{notes/ece1228-electromagnetic-theory/}
%\newcommand{\dateintitle}{}
%\newcommand{\keywords}{}

\input{../latex/peeter_prologue_print2.tex}

\usepackage{peeters_layout_exercise}
\usepackage{peeters_braket}
\usepackage{peeters_figures}
\usepackage{siunitx}
\usepackage{verbatim}
%\usepackage{mhchem} % \ce{}
%\usepackage{macros_bm} % \bcM
%\usepackage{macros_qed} % \qedmarker
%\usepackage{txfonts} % \ointclockwise

\beginArtNoToc

\generatetitle{Newton's method for reciprocal}
%\chapter{Newton's method for reciprocal}
%\label{chap:newtons_reciprocal}
% \citep{sakurai2014modern} pr X.Y
% \citep{pozar2009microwave}
% \citep{qftLectureNotes}
% \citep{doran2003gap}
% \citep{jackson1975cew}
% \citep{griffiths1999introduction}

Define

\begin{equation*}
f(x) = \inv{x} - b,
\end{equation*}
for which
\begin{equation*}
f'(x) = -\frac{1}{x}
\end{equation*}

The slope at some initial guess \( x_1 \) is
\begin{equation*}
\begin{aligned}
f'(x_1)
&= - 1/x_1^2 \\
&\approx \frac{f(x_1) - 0}{x_1 - x_2}.
\end{aligned}
\end{equation*}

Rearranging for \( x_2 \), we have
\begin{equation*}
x_2 = x_1 + x_1 (1 - b x_1)
\end{equation*}

%}
%\EndArticle
\EndNoBibArticle
